\section{Related Work}
\label{sec:related}

\subsection{Global Path Planning}
Classical graph-search algorithms provide the routing foundation for many mobile-robot systems. Dijkstra's algorithm gives shortest paths on graphs with nonnegative edge weights \cite{dijkstra1959note}, while A* accelerates search by adding an admissible heuristic \cite{hart1968formal}. For grid maps, variants such as Jump Point Search exploit path symmetries \cite{harabor2011online}, and Theta* relaxes grid-constrained headings to produce any-angle paths \cite{nash2007theta}. Recent surveys continue to emphasize that path planning must balance optimality, efficiency, map representation, and robustness in mobile-robot applications \cite{liu2023pathplanning}. These methods are effective for connecting navigation states, but they do not decide which navigation states should be visited for manipulation. This paper therefore treats global planners as replaceable routing modules after reachability-aware docking poses have been selected.

\subsection{Mobile Manipulation and Base Placement}
Mobile manipulation planning combines base motion with arm feasibility, grasping, and task execution \cite{siciliano2009robotics,lynch2017modern}. Recent reviews and task-and-motion planning studies further show that mobile manipulation requires coordination between symbolic task decisions, robot placement, perception, and continuous motion feasibility \cite{sandakalum2022motion,zhao2024survey,zhang2022grop}. Learning-based systems and benchmarks have extended this setting to object rearrangement, bimanual mobile manipulation, last-mile navigation, and tabletop rearrangement \cite{gu2023multiskill,fu2024mobilealoha,zhang2025momakitchen,hu2025mobiletabletop}. These works motivate the need to reason about where the mobile base should stand before manipulation is attempted. The present paper focuses on a complementary geometric planning layer: given tabletop task points and known scene geometry, it selects one or more reachable docking poses per table and then connects them with a global route.

\subsection{Reachability-Aware Docking and Base-Pose Selection}
A common strategy in mobile manipulation is to sample feasible base placements and evaluate manipulation reachability or inverse-kinematic feasibility. Recent work has combined navigation and manipulation costs for robot placement \cite{reister2022combining}, learned reachability behavior priors \cite{jauhri2022reachability}, considered approach planning under pose uncertainty \cite{yamazaki2022approaching}, optimized base placement for coverage mobile manipulation \cite{zhang2023baseplacement}, and used manipulability or dexterity maps to characterize useful arm workspaces \cite{sereinig2024comfort,yao2024edm}. Base-pose sequence planning has also been studied for pickup tasks \cite{naik2024basenet}. These studies are closely related to the docking-pose candidates used here. The difference is that the present problem involves many tabletop task points across many tables, so a base placement is valuable only through the subset of tasks it covers and through its effect on the global route. The proposed framework therefore treats base placement as a constrained coverage decision rather than an isolated local feasibility test.

\subsection{Coverage, Viewpoint, and Service-Pose Selection}
Coverage planning studies how a robot can visit or observe regions of interest \cite{choset2001coverage}. Viewpoint planning and frontier-based exploration similarly select poses that cover sensing targets or unknown space \cite{yamauchi1997frontier}. These problems share a set-cover structure with docking-pose selection: each candidate pose covers a subset of targets, and the planner must choose a useful subset. However, tabletop manipulation differs from sensing coverage because coverage depends on dual-arm reachability, table-side geometry, safety clearance, and docking orientation. The candidate score in this paper augments coverage with these manipulation-specific constraints.

\subsection{Hierarchical and Constraint-Based Planning}
Hierarchical planning reduces a large coupled planning problem into structured subproblems. In mobile manipulation, this separation is attractive because full joint optimization over task assignment, base route, and arm motion can be computationally expensive \cite{zhao2024survey}. Set cover and greedy selection provide practical approximations for selecting service poses \cite{chvatal1979greedy}, while routing heuristics are often used after goal states are chosen \cite{lavalle2006planning}. The proposed framework follows this hierarchical philosophy but makes the interface explicit: local candidates expose coverage and reachability sets, and the global layer only routes selected docking poses. This separation enables the same docking decisions to be tested with A*, Dijkstra, JPS, GBFS, Theta*, Lazy Theta*, and Convex A* planners in the experiments.
